\section{Criptografía de Clave Pública Clásica y su Vulnerabilidad Cuántica}
% =============================================================================

Los esquemas criptográficos de clave pública empleados históricamente de forma predominante en sistemas de votación electrónica ---entre ellos RSA y ECDSA--- sustentan su seguridad computacional en dos problemas aritméticos fundamentales: la factorización de enteros grandes y el logaritmo discreto sobre grupos finitos o curvas elípticas.

\subsection{El Mecanismo Clásico de RSA}

Propuesto por Rivest, Shamir y Adleman (1978), el esquema RSA basa su asimetría en que la multiplicación de dos números primos grandes $p$ y $q$ es trivialmente calculable en milisegundos, mientras que la operación inversa ---factorizar el módulo compuesto $N = p \cdot q$--- requiere un esfuerzo computacional que crece de forma subexponencial con el algoritmo clásico más avanzado conocido (la Criba General del Cuerpo de Números, GNFS).

La función de cifrado de un mensaje $m \in \mathbb{Z}_N$ bajo el exponente público $e$ se define como:
\begin{equation}
	c \equiv m^e \pmod N
\end{equation}
Para descifrar, el receptor legítimo requiere calcular el exponente privado $d \equiv e^{-1} \pmod{\phi(N)}$, donde $\phi(N) = (p-1)(q-1)$ corresponde a la función totiente de Euler:
\begin{equation}
	m \equiv c^d \pmod N
\end{equation}
Un atacante que intente deducir $d$ únicamente a partir de la información pública $(e, N)$ se enfrenta a la necesidad insalvable de factorizar $N$. Para una clave estándar de RSA-2048, el esfuerzo computacional clásico demandaría miles de años de cómputo en supercomputadoras actuales.

\subsection{El Algoritmo Cuántico de Shor}

La seguridad de RSA y de los esquemas basados en curvas elípticas colapsa teóricamente frente al advenimiento de una computadora cuántica de escala suficiente debido al algoritmo de factorización y logaritmo discreto formulado por Peter Shor (1997).

Shor demostró que el problema de factorizar un entero compuesto $N$ puede reducirse matemáticamente al problema de encontrar el \textit{período} $r$ de una función periódica definida sobre el grupo multiplicativo $(\mathbb{Z}_N)^*$:
\begin{equation}
	f(x) = a^x \pmod N
\end{equation}
donde $a$ es un entero coprimo con $N$. Dicha función es estrictamente periódica: satisface $f(x + r) = f(x)$, con $a^r \equiv 1 \pmod N$. Si el período $r$ es par, se cumple:
\begin{equation}
	a^r - 1 \equiv (a^{r/2} - 1)(a^{r/2} + 1) \equiv 0 \pmod N
\end{equation}
Al calcular el máximo común divisor mediante el algoritmo euclídeo clásico:
\begin{equation}
	\gcd(a^{r/2} - 1, N) \quad \text{y} \quad \gcd(a^{r/2} + 1, N)
\end{equation}
se obtienen de forma no trivial los factores primos $p$ y $q$ de $N$.

Mientras que una computadora clásica requiere tiempo exponencial o subexponencial para hallar dicho período evaluando combinaciones punto por punto, una computadora cuántica ejecuta la \textbf{Transformada Cuántica de Fourier} (\textit{Quantum Fourier Transform}, QFT). La QFT aprovecha el fenómeno de superposición cuántica y la interferencia destructiva para concentrar las amplitudes de probabilidad en las frecuencias resonantes correspondientes al período $r$, resolviendo la búsqueda en tiempo polinomial estricto $\mathcal{O}((\log N)^3)$.

\subsection{La Amenaza ``Cosechar Ahora, Descifrar Después''}

La vulnerabilidad cuántica no representa únicamente un riesgo futuro, sino una amenaza operativa inmediata clasificada por las agencias de estandarización técnica bajo el término \textit{Harvest Now, Decrypt Later} (HNDL; cosechar ahora para descifrar después) (Chen et al., 2016). Un adversario puede interceptar y almacenar pasivamente las boletas y credenciales electorales transmitidas en los procesos democráticos contemporáneos cifrados con RSA o ECDSA. Una vez que se cuente con un procesador cuántico tolerante a fallos (\textit{Fault-Tolerant Quantum Computer}, FTQC), dicho registro histórico podrá ser descifrado retroactivamente, vulnerando permanentemente la confidencialidad del voto de los electores.

% =============================================================================